\section{Introduction}\label{sec:intro}

Britain has just lived through the densest sequence of trade shocks in
its post-war history: a new customs border with its largest trading
partner (2021), the sharpest terms-of-trade deterioration since the
1970s (2022), a tariff war with its largest single-country export
market (2025), and two new trade agreements (2024, 2026). Each episode
has its own literature; none has been studied at the level where trade
policy is actually felt---the household, through the tax--benefit
system. This paper studies all of them there, together, under one
lens.

The design holds the second stage fixed and varies the shock. Each
episode enters as a \emph{declared first stage}---published estimates
for the negative shocks, the government's own projections for the
positive ones, never re-estimated here---and is mapped onto households
through two channels: consumer prices against expenditure patterns,
and labour income through the statutory tax--benefit system. Because
gross sizes span two orders of magnitude, comparisons run in ratios:
burden as a share of spending by income decile, cushioning shares, and
per-pound-of-shock intensities. And because the rulebook itself moved
across the window---the Covid uplift removed in 2021, upratings lagging
a year behind prices, the Energy Price Guarantee legislated
mid-shock---the design treats the tax--benefit vintage as an explicit
dimension: the same shock is scored under the rules of its day, under
the previous year's rules, and with the discretionary layer switched
off.

The first-pass results already discipline the received wisdom in three
ways. First, the negative shocks were regressive and the positive ones
were not progressive: the energy shock put
\MsEnergyGrossPctSpendDecOne\ per cent of bottom-decile spending at
stake against \MsEnergyGrossPctSpendDecTen\ per cent at the top, the
TCA food-price effect \MsTcaPctSpendDecOne\ against
\MsTcaPctSpendDecTen\ per cent, while the India agreement's consumer
gains are proportional and worth pounds-per-year, not per cent.
Second, per pound of gross shock, the episodes differ in bite: the
energy shock was the most bottom-decile-intensive
(\MsEnergyNormBottomPctPerBn\ points of bottom-decile spending per
\pounds 1 billion), the FTA gains the flattest
(\MsCetaNormBottomPctPerBn). Third---the finding the rulebook axis
exists for---the British state cushioned its trade shocks with
different instruments on different margins: the earnings-led tariff
shock was absorbed automatically (\MsUSTariffDisplacementCushionPct\
to \MsUSTariffWageCutCushionPct\ per cent, from the companion study),
the price-led energy shock was absorbed by discretion (the Energy
Price Guarantee took \MsEnergyCushionSharePct\ per cent of the gross
shock), and the automatic price-side mechanism---indexation---ran a
year late, costing a Universal Credit household
\pounds\MsUcLagCostGbp\ (\MsUcLagCostPctAllowance\ per cent of the
allowance) in the crisis year, months after a
\pounds\MsUcUpliftRemovalGbp\ uplift had been withdrawn.

The contribution is threefold. First, comparability: no study, for any
country, has run multiple trade shocks through a common statutory
household lens; the closest traditions---the World Bank's tariff
incidence framework, the CGE-microsimulation linkages, the EUROMOD
shock literature---each supply one ingredient
(Section~\ref{sec:literature}). Second, the rulebook axis: holding a
shock fixed and varying the vintage of the welfare state that receives
it has, to our knowledge, no precedent, and post-Covid Britain---four
shocks landing on four rulebooks---is its natural laboratory. Third,
honesty about epistemic status as a design principle: every first
stage is labelled estimate or projection, every dial is declared, and
one episode (CPTPP) is included purely as a placebo to demonstrate
that the framework does not manufacture distributional findings from
near-zero shocks.

The paper is a working draft in the same staged convention as its
companion \citep{ahmadi2026}: the first pass runs the declared first
stages against published expenditure-by-decile tables and statutory
arithmetic (Sections~\ref{sec:method}--\ref{sec:results}); the
household-record stage---PolicyEngine UK on the enhanced Family
Resources Survey, whose consumption imputation and parameter history
have been verified for exactly this use---is specified in
Section~\ref{sec:roadmap}. Section~\ref{sec:framework} defines the
estimand, Section~\ref{sec:episodes} the five episodes and their
declared first stages, and Section~\ref{sec:discussion} concludes.
